\chapter{Experiment 1: Traceability Evaluation (RQ1, RQ2)}
\label{chap:experiment-traceability}

This chapter reports the traceability evaluation for RQ1 and RQ2. It first
explains why semantic support is more suitable than word similarity for this
task. It then compares the three single-pass strategies and evaluates the
proposed \emph{spec-guided} strategy. Finally, it examines the low-scoring
requirements to clarify what the automatic results actually represent.

Experiment~1 uses two project histories and three AI coding agents. For each
agent and project, four master-prompt construction strategies are evaluated:
\emph{basic}, \emph{loose}, \emph{strict}, and \emph{spec-guided}. The first
three strategies produce the eighteen single-pass candidates used to answer
RQ1. The six \emph{spec-guided} candidates are used to answer RQ2, giving 24
candidates in total. Each candidate is evaluated at the atomic-requirement
level using mean AlignScore and untraceable rate. The complete experimental
procedure is described in \cref{chap:methodology}, while the exact project
histories and instructions are reproduced in \cref{app:datasets}.

\section{Why Word Similarity Is Not Enough}
\label{sec:demo-experiment}

A faithful master prompt must preserve the meaning of the original
requirements. Similar wording alone is not sufficient because changing one
value, condition, or permission can reverse the meaning of a requirement.

Before AlignScore was selected, it was compared with BERTScore precision
\citep{zhang2020bertscore} on two sentence pairs. The first pair contains
almost identical wording but opposite meanings. The second is a faithful
paraphrase taken from a real candidate.

\begin{lstlisting}[breaklines=true, basicstyle=\small\ttfamily, caption={Two sentence pairs scored with BERTScore precision and AlignScore.}, label={lst:demo-pairs}]
Pair 1 -- contradiction
Ground truth: "A navigation label starting with H must have a blue background."
Candidate:    "A navigation label starting with H must have a red background."
BERTScore (P) = 0.98        AlignScore = 0.0002

Pair 2 -- faithful paraphrase
Ground truth: "Implement customer management with create, edit, delete, and browse operations."
Candidate:    "Include customer management: Users can create, edit, delete, and browse customers."
BERTScore (P) = 0.63        AlignScore = 0.87
\end{lstlisting}

In Pair~1, only the color changes, but this makes the candidate contradict the
ground truth. BERTScore still gives it $0.98$ because most of the wording is
unchanged. AlignScore instead checks whether the candidate is supported by the
ground truth, so the score falls to $0.0002$.

Pair~2 expresses the same customer-management requirement with different
wording. BERTScore gives it only $0.63$, while AlignScore recognizes the
semantic support and gives it $0.87$.

These examples show why an entailment-based metric is used in this thesis.
Traceability requires distinguishing a contradiction from a faithful
paraphrase, not only measuring how similar two sentences appear.

\section{Results on the Single-Pass Strategies}
\label{sec:results-18}

RQ1 compares three single-pass strategies: \emph{basic}, \emph{loose}, and
\emph{strict}. Each strategy was tested with three agents on two projects,
producing eighteen master-prompt candidates.

Each candidate was divided into atomic requirements using the procedure
described in \cref{sec:method-metrics}. Every candidate requirement was then
scored against the complete ground-truth master prompt for the same project.

Two values are reported:

\begin{itemize}[leftmargin=*]
  \item \textbf{Mean AlignScore}: the average score across all requirements
  in the candidate;
  \item \textbf{Untraceable rate}: the percentage of requirements with an
  AlignScore below $0.5$.
\end{itemize}

A requirement below $0.5$ is treated as requiring further inspection. It is
not automatically considered non-traceable because some faithful requirements also
receive unexpectedly low scores.

\begin{table}[htbp]
  \centering
  \caption{Requirement-level AlignScore results for the simple project, reported as mean AlignScore (untraceable rate).}
  \label{tab:results-simple}
  \begin{tabular}{@{}lcccc@{}}
    \toprule
    Agent & \emph{basic} & \emph{loose} & \emph{strict} & \emph{spec-guided} \\
    \midrule
    Claude & 0.086 (92\%) & 0.824 (10\%) & 0.988 (0\%) & 0.989 (0\%) \\
    Codex  & 0.387 (62\%) & 0.857 (18\%) & 0.828 (18\%) & 0.987 (0\%) \\
    Gemini & 0.083 (97\%) & 0.737 (21\%) & 0.917 (10\%) & 0.989 (0\%) \\
    \bottomrule
  \end{tabular}
\end{table}

\begin{table}[htbp]
  \centering
  \caption{Requirement-level AlignScore results for the hard project, reported as mean AlignScore (untraceable rate).}
  \label{tab:results-hard}
  \begin{tabular}{@{}lcccc@{}}
    \toprule
    Agent & \emph{basic} & \emph{loose} & \emph{strict} & \emph{spec-guided} \\
    \midrule
    Claude & 0.382 (62\%) & 0.362 (67\%) & 0.905 (3\%) & 0.882 (11\%) \\
    Codex  & 0.787 (19\%) & 0.772 (21\%) & 0.860 (10\%) & 0.895 (8\%) \\
    Gemini & 0.391 (68\%) & 0.414 (60\%) & 0.865 (10\%) & 0.785 (17\%) \\
    \bottomrule
  \end{tabular}
\end{table}

The clearest difference is between \emph{basic} and \emph{strict}. The
\emph{strict} strategy achieves a higher mean AlignScore and a lower
untraceable rate for all six agent--project pairs. For example, Claude rises
from $0.086$ to $0.988$ on the simple project, while Gemini rises from
$0.391$ to $0.865$ on the hard project.

The improvement is not steady from \emph{basic} to \emph{loose} to
\emph{strict} in every case. On the hard project, \emph{loose} performs
slightly worse than \emph{basic} for Claude and Codex. On the simple project,
Codex has a slightly higher mean score with \emph{loose} than with
\emph{strict}, while both have the same untraceable rate.

Nevertheless, \emph{strict} outperforms \emph{loose} in five of the six
comparisons. The results therefore support the use of an explicit
traceability rule. They do not support the broader claim that adding more
instructions always leads to steady improvement.

These scores evaluate only the requirements included in a candidate. They do
not show whether the candidate omitted a valid ground-truth requirement.
Ground-truth coverage is therefore not measured in this experiment.

\section{The Spec-Guided Strategy}
\label{sec:rq2-design}

The three RQ1 strategies create the master prompt once, after the complete
prompt history has been provided. RQ2 examines whether traceability improves
when the master prompt is maintained throughout the session.

The proposed \emph{spec-guided} strategy maintains two files:

\begin{itemize}[leftmargin=*]
  \item \texttt{prompts.md}, an append-only record of the original prompts;
  \item \texttt{spec.md}, the current master prompt containing the
  requirements that remain valid.
\end{itemize}

After each new prompt, the instruction is recorded and \texttt{spec.md} is
updated. Compatible requirements are added. Clear modifications, removals,
and undo instructions update the affected topic. If a new instruction
conflicts with a settled requirement without clearly revising it, the conflict
is reported rather than silently resolved.

The complete procedure is described in \cref{sec:method-pipeline}, and the
governance instruction is reproduced in \cref{lst:governance}. During its
development, two recurring problems were found.

\section{Developing the Governance Instruction}
\label{sec:rq2-governance}

Both problems came from the same behavior: the agents tended to edit the
existing sentence using the latest prompt's wording instead of writing a
complete requirement from the relevant prompt history.

\subsection{Failure Mode 1: Edit Instructions Remain in the Specification}
\label{sec:rq2-edit-wording}

Some prompts describe how an earlier requirement should change rather than
directly stating the final requirement. For example:

\begin{quote}
  \emph{Use top navigation instead.}
\end{quote}

and:

\begin{quote}
  \emph{Change navigation buttons: Dashboard, Policies, History, Contact.}
\end{quote}

These instructions are clear in the original conversation because the earlier
state is still available. In a standalone master prompt, however, ``instead''
refers to information that may no longer be present, while ``Change
navigation buttons'' describes an edit rather than the final requirement. 

Early versions of the governance preserved this wording in
\texttt{spec.md}. Simply removing individual words such as ``instead'' or ``change''
was not a reliable solution. For example, deleting ``Change'' left
the requirement without a verb.

The final rule therefore works at the topic level. When a new prompt affects
an existing topic, the agent must read the relevant prompts and write one new
sentence stating the current requirement. It must not patch the previous
sentence with the latest wording.

For example:

\begin{quote}
  \emph{Change navigation buttons: Dashboard, Policies, History, Contact.}
\end{quote}

becomes:

\begin{quote}
  \emph{Add 4 navigation buttons: Dashboard, Policies, History, Contact.}
\end{quote}

This rule also applies to changed names, values, counts, and conditions.

\subsection{Failure Mode 2: Conversational Wording Remains in the Specification}
\label{sec:rq2-conversational-wording}

The hard project contains more conversational wording. For example:

\begin{quote}
  \emph{I'd like to build a simple ERP web application for a small business.
  Let's start with a clean foundation that we can gradually expand.}
\end{quote}

The actual requirements are to build the application and give it a clean,
expandable foundation. Expressions such as ``I'd like to'' and ``Let's start
with'' belong to the conversation rather than to the final specification.

Similar cases include temporary comments or reasons for requesting a feature,
such as:

\begin{quote}
  \emph{The dashboard doesn't need much yet.}
\end{quote}

and:

\begin{quote}
  \emph{Since we're buying from suppliers, let's add vendor management too.}
\end{quote}

These expressions originate from real prompts, but not all of their wording
describes what the system must implement.

The same topic-level rule also addresses this problem. The agent keeps the
requested names, values, conditions, and constraints, but writes them as a
standalone requirement.

For example:

\begin{quote}
  \emph{Build a simple ERP web application for a small business with a clean
  foundation that can be gradually expanded.}
\end{quote}

The two failure modes therefore share one solution: reconstruct the current
requirement from the relevant prompt history instead of preserving the latest
prompt as a local sentence edit.

\begin{lstlisting}[breaklines=true, basicstyle=\small\ttfamily, caption={Examples observed while developing the spec-guided governance instruction.}, label={lst:rq2-governance-examples}]
simple project

Navigation buttons
Before: "Change navigation buttons: Dashboard, Policies, History, Contact."
After:  "Add 4 navigation buttons: Dashboard, Policies, History, Contact."

Navigation layout
Before: "Use top navigation instead."
After:  "Use top navigation."

hard project

Project introduction
Before: "I'd like to build a simple ERP web application for a small business. Let's start with a clean foundation that we can gradually expand."
After:  "Build a simple ERP web application for a small business with a clean foundation that can be gradually expanded."
\end{lstlisting}

\section{RQ2 Results}
\label{sec:rq2-results}

The \emph{spec-guided} strategy performs better than the same agent's
\emph{strict} result in four of the six comparisons.

\begin{table}[htbp]
  \centering
  \caption{\emph{Spec-guided} compared with the same agent's \emph{strict} result. Bold marks the better result in each row.}
  \label{tab:rq2-vs-strict}
  \begin{tabular}{@{}llcc@{}}
    \toprule
    Project & Agent & \emph{strict} & \emph{spec-guided} \\
    \midrule
    simple project & Claude & 0.988 (0\%)  & \textbf{0.989} (\textbf{0\%})  \\
    simple project & Codex  & 0.828 (18\%) & \textbf{0.987} (\textbf{0\%}) \\
    simple project & Gemini & 0.917 (10\%) & \textbf{0.989} (\textbf{0\%})  \\
    hard project   & Claude & \textbf{0.905} (\textbf{3\%})  & 0.882 (11\%)  \\
    hard project   & Codex  & 0.860 (10\%) & \textbf{0.895} (\textbf{8\%})  \\
    hard project   & Gemini & \textbf{0.865} (\textbf{10\%}) & 0.785 (17\%) \\
    \bottomrule
  \end{tabular}
\end{table}

On the simple project, all three \emph{spec-guided} candidates reach a
$0\%$ untraceable rate. The final outputs also avoid the edit-oriented
navigation wording found during development.

On the hard project, only Codex improves on its \emph{strict} result. Its mean
AlignScore rises from $0.860$ to $0.895$, while its untraceable rate falls
from $10\%$ to $8\%$. Claude and Gemini both receive lower automatic
scores with \emph{spec-guided}.

However, the low-scoring requirements do not all represent invented content.
Manual review identified three different causes.

\subsection{Review of Low-Scoring Requirements}

\textbf{Same requirements with inconsistent wording.}
Some low-scoring candidate requirements describe the same system behavior as
the ground truth but use shorter or less consistent requirement wording. For
example:

\begin{quote}
\emph{Introduce user roles: Admin, Staff, and Manager}
\end{quote}

receives an AlignScore of $0.265$ against:

\begin{quote}
\emph{Support the user roles Admin, Manager, and Staff.}
\end{quote}

Similarly, \emph{Add sales orders} receives a score of $0.274$ against
\emph{Implement sales orders}. In both cases, the candidate preserves the
intended feature, but uses a different requirement verb from the ground truth.

These examples show that, although the governance is designed to keep the
output close to the original prompts, it does not always produce consistent
software-requirement wording. The agents may use verbs such as ``Add'' or
``Introduce'' where the ground truth uses ``Implement'' or ``Support.''
AlignScore also has difficulty recognizing that these sentences express the
same software requirement despite their different wording. The same low
scores were obtained when the sentence pairs were evaluated separately,
showing that the results were repeatable rather than caused by the other
requirements in the candidate.

\textbf{Non-requirement content.}
Some low-scoring text comes directly from the prompts but expresses a reason,
transition, or temporary comment rather than a requirement that must remain in
the master prompt.

For example, the request for an Excel export explains that the file should be
shared with an accountant. The export feature is a requirement, while the
accountant clause is the reason for requesting it.

Expressions such as \emph{The dashboard doesn't need much yet} create the
same problem. They are not invented, but they do not state a complete
requirement.

\textbf{Hallucinated details.}
Gemini's hard-project candidate contains two clear additions that do not
appear in the original prompts:

\begin{itemize}[leftmargin=*]
  \item an Admin-only User Accounts management view;
  \item Vendors added to a navigation bar that was requested to contain only
  Dashboard, Customers, and Products.
\end{itemize}

These details are hallucinations as they introduce system behavior or
interface elements that cannot be traced to the original prompt history. They
show that repeatedly applying the no-invention rule reduces, but does not
completely prevent, hallucinated content from entering the maintained
specification.

\subsection{Interpretation}

The manual review shows that the automatic results should not be read as a
direct count of hallucinated requirements. The requirements below the $0.5$
threshold fall into the three categories described above: same requirements
with inconsistent wording, non-requirement content, and hallucinated
details.

Codex shows the clearest improvement on the hard project. Its mean AlignScore
increases from $0.860$ to $0.895$, while its untraceable rate decreases from
$10\%$ to $8\%$. Most of its low-scoring units either describe the same
system behavior as the ground truth but use different requirement wording, or
contain incomplete conversational expressions. The manual review did not find
comparable hallucinated functionality in these units. The automatic and manual
results therefore both indicate an improvement over its \emph{strict}
baseline.

Claude's untraceable rate increases from $3\%$ to $11\%$, but this does not
mean that $11\%$ of its requirements are hallucinations. Several flagged
requirements describe the same system behavior as the ground truth but use
shorter or different verbs, such as ``Add'' instead of ``Implement'' or
``Introduce'' instead of ``Support.'' Although the governance is designed to
keep the output close to the original prompts, it does not always produce
consistent software-requirement wording. AlignScore also has difficulty
recognizing that these sentences express the same software requirement
despite their different wording. Claude additionally retains some unnecessary
conversational wording, but no hallucinated feature was found during the
manual review. Its result is therefore closer to the \emph{strict} baseline
than the raw untraceable rate suggests.

Gemini shows a different result. Some of its flagged requirements also use
inconsistent wording or preserve conversational material, but its candidate
contains two hallucinated details: an Admin-only User Accounts management view
and Vendors added to the navigation bar. Its lower score therefore reflects a
real decrease in traceability. However, the $17\%$ untraceable rate should not
be interpreted as the percentage of hallucinated requirements because only
part of that rate comes from the two hallucinations.

The untraceable rate is therefore most useful as a screening measure. It shows
which requirements require manual inspection, but it does not identify the
cause of each low score. A flagged unit may express the same requirement with
inconsistent wording, retain text from the conversation that is not a final
requirement, or introduce content that was never requested.

Overall, the \emph{spec-guided} strategy improves on \emph{strict} in four
of the six agent--project comparisons. It performs strongly for all three
agents on the shorter simple project, but its results are mixed on the longer
hard project: Codex improves, Claude remains close to its baseline after
manual review, and Gemini introduces two hallucinated details.

RQ2 therefore does not show that per-turn consolidation is always more
traceable than a single final consolidation. Maintaining \texttt{spec.md}
reduces the amount of history that must be reconstructed in one step, but
every update requires the agent to interpret the new prompt and rewrite the
affected requirement correctly. As the history becomes longer, these repeated
updates create more opportunities for inconsistent wording, unnecessary
conversational text, or hallucinated details to enter the specification. This
trade-off is discussed further in \cref{sec:discussion-tradeoff}.

